\documentclass[12pt]{article}
\usepackage[utf8]{inputenc}
\usepackage[T1]{fontenc}
\usepackage{hyperref}
\title{Xicotepec de Ju\'arez Totonac Language Resources for K-12 Curriculum Development}
\author{Compiled by Claude Sonnet 5 (claude-sonnet-5) for Andr\'as Kornai}
\date{}
\begin{document}
\maketitle

\section{Introduction}
Xicotepec de Ju\'arez Totonac (ISO 639-3 \texttt{too}; Glottocode xico1235) is a Totonacan language of Puebla, Mexico, with about 3,000 speakers according to Endangered Languages Project pages citing Ethnologue (2009). Sources disagree on how endangered it is, and it is not known to be taught in schools. Few resources turned up in the search, and most are old.

\section{Primary Educational Resources}
\begin{itemize}
\item \emph{Lecturas en Totonaco} (1981), a primer and reader in Totonac and Spanish held by El Colegio de M\'exico. It describes the alphabet as mostly Spanish letters plus a few extra ones. \url{https://repositorio.colmex.mx/concern/books/hx11xh02x?locale=en}
\item A second item in the same repository (a vocabulary, 1965). Its record notes that its language is no longer current. \url{https://repositorio.colmex.mx/downloads/zc77st91s?locale=es}
\end{itemize}

\section{Linguistic and Reference Materials}
\begin{itemize}
\item \emph{Diccionario totonaco de Xicotepec de Ju\'arez, Puebla} (Instituto Ling\"u\'istico de Verano / SEP). I saw it only as a citation in the ASJP entry, not as a downloadable file. \url{https://asjp.clld.org/languages/XICOTEPEC_TOTONAC}
\item Glottolog language page, with classification and bibliography. \url{https://glottolog.org/resource/languoid/iso/too}
\item Endangered Languages Project entry for Northern Totonac. \url{https://endangeredlanguages.com/node/38948}
\item Wikitongues language page. \url{https://wikitongues.org/languages/too/}
\end{itemize}

\section{Religious and Community Resources}
\begin{itemize}
\item Scripture Earth lists a New Testament, Old Testament selections and an Android app. \url{https://scriptureearth.org/00eng.php?iso=too}
\end{itemize}
I found no community organisation or school programme in the search results.

\section{Digital and Audio Resources}
\begin{itemize}
\item Global Recordings Network offers evangelistic and Bible-teaching audio. \url{https://globalrecordings.net/en/language/too}
\end{itemize}

\section{Curriculum Development Context}
Material found so far is limited to a 1981 primer, an older dictionary, scripture and audio. Any curriculum would need new native-language materials developed with the community. The scripture and dictionary texts can serve as a source of orthography and vocabulary. I did not verify the licensing or current availability of any item.

\end{document}